\begin{abstract}\noindent
\textbf{Purpose} -- To formulate, analyse and partly verify a conjugate heat-transfer problem in which a thin
polymer strip, heated on its outer face, forms one wall of a laminar slit channel and bows because of through-wall
thermal expansion, widening the passage.\\
\textbf{Design/methodology/approach} -- A steady formulation with one consistent set of sign and interface
conventions is given. Closed-form thermoelastic results (free curvature, bowing, buckling, beam--column
amplification) define a baseline and admissibility guards. A finite-element method and a
fixed-point coupling algorithm are specified; the wall is verified with two independent plane-strain codes, including a restrained,
geometrically nonlinear case.\\
\textbf{Findings} -- For water at $\mathrm{Re}=\NumRe$ and a flux of \SI{5}{\kilo\watt\per\square\metre} on a
\SI{3}{\milli\metre} PEEK-class strip, the strip bows away from the channel by $\NumWscH\,H_0$
(reduced order) or $\NLwFEH\,H_0$ (finite element); the difference of $\NLoverScPct\,\%$ vanishes for slender strips. Deformation and stress,
not buckling, bound the flux. Indicative estimates suggest a roughly halved pressure drop and a higher wall excess
temperature, while the return effect of the flow on the bow is small.\\
\textbf{Research limitations/implications} -- No coupled fluid--structure result is reported; the flow-side
estimates are reduced-order, the material data are representative; contact, heater and casing are idealised.\\
\textbf{Originality/value} -- Thermally driven wall deformation in forced convection, a verified wall model and a specified, not yet executed, coupled solver.\\[0.5ex]
\textbf{Keywords} Conjugate heat transfer, Thermoelasticity, Fluid--structure interaction, Finite element method,
Flexible channel wall, Thermal bending\\
\textbf{Paper type} Research paper
\end{abstract}
